\subsection{5 安全红队测试}\label{safety-red-teaming}

为了评估自动化基准之外的安全特性，我们与模型开发周期并行开展了红队测试。目标是揭示自动化评估无法可靠检测的对抗性漏洞、新型攻击向量和危害类别缺口，并将发现结果按滚动节奏反馈到训练数据收集与策略细化中。

\subsection{5.1 内部红队测试}\label{internal-red-teaming}

方法论。内部红队测试由 MAI 红队成员（安全研究人员和招募的外部标注人员）执行，并在整个开发周期中对不同模型版本开展。



图 20. 按危害类别的安全--有用性权衡。指标为每个提示两次生成结果的平均值。对于给定的危害类别，安全--有用性权衡更好的模型位于另一模型的右上方。两个模型之间用虚线连接。

\pandocbounded{\includegraphics[keepaspectratio,alt={image}]{images/4117e1df1c48425fa29f3357a91cff9ecb08956a230d298d448fd3ac270d06f4.jpg}}

在覆盖 MAI-Thinking-1 开发早期、中期和后期阶段的 15 次合作中，红队成员在 25 个政策类别上执行了超过 2,170 个基于目标的对抗性场景。每个场景运行 5--10 轮对话，以允许升级到首轮拒绝之外。随着后续模型版本能力增强，我们提高了评估的范围和难度。

标注人员经过对抗性提示能力筛选，并获准访问内部模型部署。他们收到一份预填充的基于目标的场景追踪表，每个场景都附有支持性上下文和目标政策类别，并被指示在多轮对话中诱导政策违规。违规情况标注了结果、严重程度、首次违规的轮次、最严重违规的文本以及其他元数据。该计划有几个范围约束：测试主要使用英语，非英语输入仅在有限情况下作为越狱向量使用，而非系统性评估；智能体工具使用和多模态输入不在范围内。

发现与缓解措施。该计划的一个核心产出是攻击模式分类法，这些模式在红队成员和模型检查点之间独立反复出现，我们认为这证明这些模式（而非单个提示）才是持久的对抗面。其中六种模式最为突出：在良性借口下进行多轮升级；虚构或小说式框架；凭证化身份借口；逐步递归或格式漂移（反复要求扩展、重排格式或落实先前含糊其辞的回答）；上下文内年龄指示绕过；以及权威文档伪造。在各类别中，同样的模式驱动了大多数红队成功案例，该分类法现在作为后续周期中提示收集和评判准则开发的覆盖清单。

发现结果按滚动节奏反馈到训练中。一个突出的反复出现的模式是虚构框架绕过，会暴露自残内容，包括多语言变体。该模式催生了一组精选数据，被扩展进安全 evalulation，并用于对抗性提示收集的种子扩展，以纳入



图 21. 按变换类型分组的越狱评估攻击成功率。基础技术（Foundational Techniques）是保留有害意图的单步变换，通过越狱包装器或提示模板等简单修改实现。组合技术（Compositional Techniques）组合多种变换或结构化改写，包括基于模板的攻击，如 PyRIT (Munoz et al., 2024)、PAP 风格变换 (Zeng et al., 2024a) 以及非英语或混合语言变体。自适应技术（Adaptive Techniques）引入交互、搜索或多轮结构，包括 TAP (Mehrotra et al., 2024) 和多轮攻击 (Russinovich et al., 2024)。攻击成功率越低表示安全性越强。误差棒表示 95\% 置信区间。第三方模型结果包含服务方侧安全过滤。

\pandocbounded{\includegraphics[keepaspectratio,alt={image}]{images/009f2d7fadb48e72a572defbd3b2c6d8b9689167eebd5bb019056e214f3fe842.jpg}}

SFT 和 RL 安全配比。

在红队测试中识别出的最高优先级修复类别中，从缓解前到最终候选版本，攻击成功率总体下降了约 22\%。我们将越狱成功率降低了约 44\%，仇恨与公平问题降低了约 43\%，儿童安全问题降低了约 30\%，心理健康攻击降低了约 20\%。

\subsection{5.2 独立红队测试}\label{independent-red-teaming}

额外的红队测试由微软 AI 红队 (AIRT) 和第三方供应商独立进行。这些合作聚焦于静态评估最薄弱的风险领域：自动化对抗攻击方法、代码与网络滥用安全、心理社会与心理健康危害，以及多语言覆盖。结构化危险能力/能力提升评估 (Phuong et al., 2024; Li et al., 2024b) 不在本次发布范围内。

自适应剪枝攻击树（TAP）攻击 (Mehrotra et al., 2024) 被揭示为一个鲁棒性缺口。作为回应，我们构建了一条针对性的对抗性数据管线：首先广泛生成现实有害场景，再通过多样化的攻击变换模板集进行扩展，然后使用 TAP 风格的自适应细化，针对当前模型优化生成的提示，直到它们产生可靠的模型特定失败。这一闭环过程将外部红队的发现转化为针对性的修复数据，大幅降低了 TAP 越狱敏感性，并在相同攻击向量上使模型达到与最先进模型相当的水平。

低资源语言框架 (Deng et al., 2024) 被标记为另一个漏洞。用英语可靠拒绝的内容在约鲁巴语、泰卢固语、阿姆哈拉语、缅甸语、高棉语和马来语中被诱导出来。我们的回应是扩展安全训练数据配比，加入多语言对抗性种子，将高回报的英语攻击模式翻译并重新定向到受影响的语言。这显著缩小了目标语言上英语/非英语之间的差距，尽管长尾低资源语言的鲁棒性仍是持续投入的领域。



\subsection{6 集群环境}\label{cluster-environment}

MAI-Thinking-1 需要一个集群环境，既要足够可组合，以支持跨数据、模型、训练、评估和服务组件的快速迭代；又要足够可扩展，以便在前沿规模上测试改进；还要足够可靠，使测得的模型质量提升而非基础设施噪声成为关注焦点。

因此，我们将集群环境视为模型开发的活跃组成部分。在训练中，目标是在保持数值正确性、确定性恢复、高 MFU 和高 goodput 的同时，最大化每个墙钟日的有效 FLOPs。在推理中，目标是在保持质量、长上下文正确性、可预测延迟和部署效率的同时，最大化每秒和每瓦特的有效 token 数。两种场景下都适用相同的系统原则：拓扑很重要，内存移动很重要，静默的正确性失败不可接受，聚合吞吐只有在转化为可靠的模型进展或可靠的推理服务能力时才有用。

本节介绍该环境的三个组成部分。第 6.1 节介绍训练集群如何将原始加速器容量转化为稳定的实验基质。第 6.2 节介绍用于衡量所分配 GPU 是否产生有用学习进展的稳定性、确定性和 goodput 指标。第 6.3 节简要介绍在微软 MAIA-200 硬件 (Microsoft, 2026) 上运行 MAI-Thinking-1 得到的推理结果。
